% Generated draft: requires mathematical and source review.
\documentclass[11pt]{article}
\usepackage[margin=25mm]{geometry}
\usepackage{fontspec}
\setmainfont{Latin Modern Roman}
\usepackage{amsmath,amssymb,mathtools,unicode-math}
\setmathfont{Latin Modern Math}
\usepackage{enumitem,needspace,xurl,hyperref}
\setlength{\parindent}{0pt}
\setlength{\parskip}{5pt}
\setlength{\emergencystretch}{3em}
\newcommand{\N}{\mathbb{N}}
\newcommand{\Z}{\mathbb{Z}}
\newcommand{\Q}{\mathbb{Q}}
\newcommand{\R}{\mathbb{R}}
\newcommand{\C}{\mathbb{C}}
\newcommand{\F}{\mathbb{F}}
\newcommand{\A}{\mathbb{A}}
\newcommand{\PP}{\mathbb{P}}
\newcommand{\E}{\mathbb{E}}
\newcommand{\eps}{\varepsilon}
\newcommand{\dd}{\,\mathrm{d}}
\begin{document}
\section*{UnsolvedMath ID: 800309}
\textbf{Source problem code:} catalogue entry
\subsection{Definitions and mathematical statement}
Let \(S^1=\R/\Z\), and let \(\operatorname{Diff}^t_+(S^1)\) denote the orientation-preserving circle diffeomorphisms of class \(C^t\). For \(f\in\operatorname{Diff}^t_+(S^1)\), choose a lift \(F\colon\R\to\R\) satisfying \(F(x+1)=F(x)+1\). Its rotation number is
\[
\rho(f)=\left[\lim_{k\to\infty}\frac{F^k(x)-x}{k}\right]\in\R/\Z.
\]
The class is independent of the lift and of \(x\). Irrational rotation number means \(\rho(f)\notin\Q/\Z\). Write \(f^k\) for the \(k\)-fold iterate and \(\operatorname{id}_{S^1}\) for the identity map.

The source asks to determine the regularity values \(r>1\) for which there exists \(s\ge r\), chosen independently of \(f\) and its irrational rotation number, such that every \(C^s\) circle diffeomorphism with irrational rotation number satisfies
\[
 f^{q_n}\longrightarrow\operatorname{id}_{S^1}
 \qquad\text{in the }C^r\text{ topology as }n\to\infty,
\]
where \((q_n)\) is described in the source as the sequence of denominators of the rational approximations of the rotation number.

\textbf{Unresolved conventions.} The supplied record does not specify which rational approximations define \((q_n)\). If the intended approximations are the regular continued-fraction convergents \(p_n/q_n\) of the representative \(\alpha(f)\in(0,1)\) of \(\rho(f)\), with each fraction reduced and \(q_n>0\), then the requested set has the following conditional formulation:
\[
\left\{r\in\R:r>1,\ \exists s\in\R,\ s\ge r,\ 
\forall f\in\operatorname{Diff}^s_+(S^1),\quad
\rho(f)\notin\Q/\Z\ \Longrightarrow\ 
f^{q_n(\alpha(f))}\xrightarrow[n\to\infty]{C^r}\operatorname{id}_{S^1}
\right\}.
\]
The continued-fraction interpretation requires confirmation. The displayed real-exponent formulation also requires a specified meaning of \(C^t\) and its topology for nonintegral \(t\); the record does not state whether only integer regularities are intended. For an integer \(r=m\), convergence in \(C^m\) means uniform convergence of the maps and their derivatives through order \(m\) in smooth local coordinates. Orientation preservation is made explicit because it is the setting of the rotation number used here; it is implicit in the source wording.
\subsection{Short English statement}
Determine which regularities \(r>1\) admit a single higher regularity \(s\ge r\) such that every \(C^s\) circle diffeomorphism with irrational rotation number has iterates approaching the identity in \(C^r\) along the specified approximation denominators. The intended approximation sequence and the convention for noninteger regularity require clarification. The supplied registry labels this problem open; that label is not independently verified.
\subsection{Sources}
\begin{itemize}
\item Local UnsolvedMath registry, permanent ID \(800309\), exact source problem code \texttt{catalogue entry}, title ``Group Actions on 1-Manifolds --- Question 26.'' This draft uses the supplied complete source record, which labels the problem open and attributes it to Navas, proposed in 2017.
\item Navas, \emph{Group Actions on 1-Manifolds: A List of Very Concrete Open Questions} (2017), source item recorded as ``Question 26 (\#25).'' Supplied source URL: \url{https://arxiv.org/abs/1712.06462}. The supplied provenance records access on 2026-10-16 and extraction by \texttt{arxiv\_latex}; these are supplied provenance claims, not independently verified access or source verification.
\end{itemize}
\end{document}
